\subsection{Probability Basics}
\tsDef{2.1 (Discrete Probability Space)}{
    A \emph{discrete probability space} is determined by a sample space
    \[
        \Omega = \{\omega_1,\omega_2,\dots\}
    \]
    of elementary outcomes.
    To each elementary outcome $\omega_i$ a probability $\Pr[\omega_i]$ is assigned such that
    \[
        0 \le \Pr[\omega_i] \le 1
    \]
    and
    \[
        \sum_{\omega \in \Omega} \Pr[\omega] = 1.
    \]
    A subset $E \subseteq \Omega$ is called an \emph{event}. The probability of an event is defined by
    \[
        \Pr[E] := \sum_{\omega \in E} \Pr[\omega].
    \]
    If $E$ is an event, then
    \[
        \bar E := \Omega \setminus E
    \]
    is called the \emph{complement} of $E$.
}
\\
\tsLem{2.2 (Basic Properties of Probabilities)}{
    For events $A,B$ the following hold:
    \begin{itemize}
        \item[(1)]
              \[
                  \Pr[\varnothing] = 0,
                  \qquad
                  \Pr[\Omega] = 1.
              \]
        \item[(2)]
              \[
                  0 \le \Pr[A] \le 1.
              \]
        \item[(3)]
              \[
                  \Pr[\bar A] = 1 - \Pr[A].
              \]
        \item[(4)] If $A \subseteq B$, then
              \[
                  \Pr[A] \le \Pr[B].
              \]
    \end{itemize}
}
\tsThe{2.3 (Addition Theorem)}{
    Let $A_1,\dots,A_n$ be pairwise disjoint events, that is,
    \[
        A_i \cap A_j = \varnothing
        \quad\text{for all } i\ne j.
    \]

    Then
    \[
        \Pr\!\left[\bigcup_{i=1}^n A_i\right]
        =
        \sum_{i=1}^n \Pr[A_i].
    \]

    Analogously, for an infinite family of pairwise disjoint events $A_1,A_2,\dots$,
    \[
        \Pr\!\left[\bigcup_{i=1}^{\infty} A_i\right]
        =
        \sum_{i=1}^{\infty} \Pr[A_i].
    \]
}
\tsThe{2.5 (Inclusion-Exclusion Principle)}{
    For events $A_1,\dots,A_n$ with $n\ge2$,
    \[
        \Pr\!\left[\bigcup_{i=1}^n A_i\right]
        =
        \sum_{l=1}^n (-1)^{l+1}
        \sum_{1\le i_1<\cdots<i_l\le n}
        \Pr[A_{i_1}\cap\cdots\cap A_{i_l}].
    \]
    Equivalently,
    \begin{align*}
        \Pr\!\left[\bigcup_{i=1}^n A_i\right]
         & = \sum_{i=1}^n \Pr[A_i]                                            \\
         & - \sum_{1\le i_1<i_2\le n}\Pr[A_{i_1}\cap A_{i_2}]                 \\
         & + \sum_{1\le i_1<i_2<i_3\le n}\Pr[A_{i_1}\cap A_{i_2}\cap A_{i_3}] \\
         & - \cdots                                                           \\
         & + (-1)^{n+1}\Pr[A_1\cap\cdots\cap A_n].
    \end{align*}
}
\tsCor{2.5 (Boolean inequality, Union Bound)}{
    For events $A_1,\dots,A_n$,
    \[
        \Pr\!\left[\bigcup_{i=1}^n A_i\right]
        \le
        \sum_{i=1}^n \Pr[A_i].
    \]
    Analogously, for an infinite sequence of events $A_1,A_2,\dots$,
    \[
        \Pr\!\left[\bigcup_{i=1}^{\infty} A_i\right]
        \le
        \sum_{i=1}^{\infty} \Pr[A_i].
    \]
}
\tsThe{Laplace Principle}{
    Under the Laplace principle, all elementary outcomes have the same probability:
    \[
        \Pr[\omega] = \frac{1}{|\Omega|}
        \quad\text{for all elementary outcomes } \omega.
    \]
    Hence for any event $E$,
    \[
        \Pr[E] = \frac{|E|}{|\Omega|}.
    \]
}